% !TEX program = lualatex
\documentclass[aspectratio=169,11pt]{beamer}

\usetheme[progressbar=frametitle,numbering=fraction]{metropolis}
\usepackage{luatexja}
\usepackage[haranoaji]{luatexja-preset}
\renewcommand{\kanjifamilydefault}{\gtdefault}

\usepackage{amsmath,amssymb}
\usepackage{booktabs}
\usepackage{tabularx}
\usepackage{xcolor}
\usepackage{listings}

% カラーパレット
\definecolor{primary}{RGB}{33, 50, 75}
\definecolor{accent}{RGB}{200, 70, 48}
\definecolor{codebg}{RGB}{246, 248, 250}
\definecolor{codeframe}{RGB}{218, 222, 229}
\definecolor{codekw}{RGB}{24, 76, 156}
\definecolor{codestr}{RGB}{3, 106, 7}
\definecolor{codecom}{RGB}{110, 119, 129}

\setbeamercolor{primary}{fg=primary}
\setbeamercolor{alerted text}{fg=accent}

\lstset{
  basicstyle=\fontsize{7.5pt}{9.5pt}\selectfont\ttfamily,
  backgroundcolor=\color{codebg},
  frame=single,
  rulecolor=\color{codeframe},
  keywordstyle=\color{codekw}\bfseries,
  stringstyle=\color{codestr},
  commentstyle=\color{codecom}\itshape,
  breaklines=true,
  tabsize=2,
  showstringspaces=false,
  aboveskip=2pt,
  belowskip=2pt
}

\title{\textbf{容量制約付き施設配置問題 (CFLP)}}
\subtitle{AMPL / amplpy による厳密解法と Timefold 解の比較・検証}
\author{数理最適化演習}
\date{\today}
\institute{lecture2026}

\begin{document}

% 1. タイトル
\maketitle

% 2. 目次
\begin{frame}{アジェンダ}
  \tableofcontents[hideallsubsections]
\end{frame}

\section{問題の概要と設定}

\begin{frame}{施設配置問題（Timefold quickstarts）の概要}
  \begin{itemize}\small
    \item \textbf{問題背景}:
      \begin{itemize}\footnotesize
        \item 複数の施設候補地から最適な場所を選定（開設）し、顧客の需要を充足させる代表的組合せ最適化問題（\alert{\textbf{Capacitated Facility Location Problem: CFLP}}）。
      \end{itemize}
    \item \textbf{データセット（ロンドン: 施設 30・顧客 60）}:
      \begin{itemize}\footnotesize
        \item 施設候補数: 30箇所（各施設に固有の開設費用 $f_j$、容量 $C_j = 150$）
        \item 顧客数: 60箇所（各顧客の需要 $d_i = 15$、合計需要 = 900）
      \end{itemize}
    \item \textbf{最適化のトレードオフ}:
      \begin{itemize}\footnotesize
        \item \textbf{開設費用 (Setup Cost)}: 開設する施設数を絞るほど安くなる。
        \item \textbf{輸送・移動距離 (Distance)}: 施設を分散配置するほど顧客への距離が短くなる。
      \end{itemize}
    \item \textbf{距離計算式 (Timefold 準拠)}:
      \begin{equation*}
        dist(i, j) = \left\lceil \sqrt{(\text{lat}_i - \text{lat}_j)^2 + (\text{lon}_i - \text{lon}_j)^2} \times 111{,}000 \right\rceil \quad \text{[m]}
      \end{equation*}
  \end{itemize}
\end{frame}

\section{記号と数理定式化}

\begin{frame}{記号の定義（Notation）}
  \begin{columns}[T]
    \begin{column}{0.48\textwidth}
      \textbf{集合 (Sets)}
      \begin{itemize}\small
        \item $I = \{1, \dots, |I|\}$: 顧客集合 ($|I|=60$)
        \item $J = \{1, \dots, |J|\}$: 施設候補集合 ($|J|=30$)
      \end{itemize}
      \vspace{0.4em}
      \textbf{重み係数 (Weights)}
      \begin{itemize}\small
        \item $w_{\text{setup}} = 2$: 施設開設費用の重み
        \item $w_{\text{dist}} = 6$: 移動距離の重み
      \end{itemize}
    \end{column}
    \begin{column}{0.5\textwidth}
      \textbf{パラメータ (Parameters)}
      \begin{itemize}\small
        \item $d_i$: 顧客 $i$ の需要量 ($=15$)
        \item $C_j$: 施設 $j$ の受入容量 ($=150$)
        \item $f_j$: 施設 $j$ の固定開設費用
        \item $dist_{ij}$: 顧客 $i$ と施設 $j$ 間の距離 (m)
      \end{itemize}
      \vspace{0.4em}
      \textbf{決定変数 (Decision Variables)}
      \begin{itemize}\small
        \item $y_j \in \{0, 1\}$: 施設 $j$ を開設するとき 1, 他は 0
        \item $x_{ij} \in \{0, 1\}$: 顧客 $i$ を施設 $j$ に割当てる時 1, 他は 0
      \end{itemize}
    \end{column}
  \end{columns}
\end{frame}

\begin{frame}{数理定式化（0-1整数線形計画問題）}
  \vspace{-0.3em}
  {\small\textbf{目的関数：加重総ペナルティの最小化}}
  \vspace{-0.3em}
  \begin{equation*}
    \min \quad \sum_{j \in J} w_{\text{setup}} f_j y_j + \sum_{i \in I}\sum_{j \in J} w_{\text{dist}} dist_{ij} x_{ij}
  \end{equation*}
  \vspace{-0.5em}
  {\small\textbf{制約条件}}
  \vspace{-0.5em}
  \begin{align}
    \sum_{j \in J} x_{ij} &= 1 & \forall i \in I & \quad \text{(単一割当制約)} \label{eq:assign} \\
    \sum_{i \in I} d_i x_{ij} &\le C_j y_j & \forall j \in J & \quad \text{(施設容量制約)} \label{eq:cap} \\
    x_{ij} &\le y_j & \forall i \in I, \forall j \in J & \quad \text{(開設連動制約 / 強定式化)} \label{eq:link} \\
    y_j &\in \{0, 1\}, \quad x_{ij} \in \{0, 1\} & \forall i \in I, \forall j \in J & \quad \text{(0-1整数変数)} \nonumber
  \end{align}
\end{frame}

\begin{frame}{各制約の数理的・実務的意味}
  \begin{itemize}\small
    \item \textbf{単一割当制約 \eqref{eq:assign}}:
      \begin{itemize}\footnotesize
        \item 各顧客 $i$ の需要はちょうど1つの施設から供給される（分割供給を禁止）。
      \end{itemize}
    \item \textbf{施設容量制約 \eqref{eq:cap}}:
      \begin{itemize}\footnotesize
        \item 開設された施設 ($y_j = 1$) の需要合計は最大容量 $C_j$ 以下。
        \item 未開設の施設 ($y_j = 0$) には需要を一切割り当てられない。
      \end{itemize}
    \item \textbf{開設連動制約 (Linking Constraint) \eqref{eq:link}}:
      \begin{itemize}\footnotesize
        \item 理論的には制約 \eqref{eq:cap} から導出可能だが、加えることで\alert{\textbf{LP緩和（線形緩和）の実行可能領域を大幅に引き締める}}（Strong Formulation）。
        \item 分枝限定法（Branch and Bound）における下界値が格段に良くなり、探索ノード数を大幅に削減。
      \end{itemize}
  \end{itemize}
\end{frame}

\section{AMPL コードの解説}

\begin{frame}[fragile]{AMPL モデル定義（Model Formulation）}
\begin{lstlisting}[language=bash]
set CUSTOMERS;   # 顧客集合 I
set FACILITIES;  # 施設候補集合 J

param demand{CUSTOMERS};          # 需要量 d_i
param capacity{FACILITIES};       # 容量 C_j
param setup_cost{FACILITIES};     # 開設費用 f_j
param dist{CUSTOMERS, FACILITIES};# 距離 dist_ij
param w_setup;                    # 開設重み (2)
param w_dist;                     # 距離重み (6)

var y{FACILITIES} binary;            # 施設開設決定変数 y_j
var x{CUSTOMERS, FACILITIES} binary; # 顧客割当決定変数 x_ij

# 目的関数：加重総コスト最小化
minimize TotalCost:
    sum{j in FACILITIES} w_setup * setup_cost[j] * y[j]
    + sum{i in CUSTOMERS, j in FACILITIES} w_dist * dist[i, j] * x[i, j];
\end{lstlisting}
\end{frame}

\begin{frame}[fragile]{AMPL 制約条件定義（Constraints）}
\begin{lstlisting}[language=bash]
# 1. 単一割当制約: 各顧客は必ず1つの施設に割り当て
s.t. AssignOnce{i in CUSTOMERS}:
    sum{j in FACILITIES} x[i, j] = 1;

# 2. 容量制約: 各施設に割り当てられた需要合計 <= 容量 * y[j]
s.t. Capacity{j in FACILITIES}:
    sum{i in CUSTOMERS} demand[i] * x[i, j] <= capacity[j] * y[j];

# 3. 強定式化連動制約: y[j] が 0 ならば x[i, j] も 0
s.t. Linking{i in CUSTOMERS, j in FACILITIES}:
    x[i, j] <= y[j];
\end{lstlisting}
  \begin{itemize}\footnotesize
    \item 数理モデルの記号表記と AMPL の構文が \alert{\textbf{1対1で直観的に対応}} している。
    \item 添字範囲（`{i in CUSTOMERS}` 等）で制約を宣言的に記述できるため、バグが混入しにくい。
  \end{itemize}
\end{frame}

\section{Python (amplpy) 実装の解説}

\begin{frame}[fragile]{Python 実装 (1)：データ読込と距離行列生成}
\begin{lstlisting}[language=Python]
import json, math
from amplpy import AMPL

def calculate_distance(loc1, loc2):
    # Timefold quickstart の距離計算式（メートル単位）
    dlat, dlon = loc1[0] - loc2[0], loc1[1] - loc2[1]
    return int(math.ceil(math.sqrt(dlat**2 + dlon**2) * 111000))

# JSONファイルから施設・顧客・重みパラメータを抽出
with open("ロンドン（施設 30・顧客 60）.json", "r", encoding="utf-8") as f:
    data = json.load(f)
problem = data.get("problem", data)
facilities, consumers = problem["facilities"], problem["consumers"]

# 距離行列の事前計算 (60 x 30 = 1,800 ペア)
dist_matrix = {
    (c["id"], f["id"]): calculate_distance(c["location"], f["location"])
    for c in consumers for f in facilities
}
\end{lstlisting}
\end{frame}

\begin{frame}[fragile]{Python 実装 (2)：AMPL へのデータ投入と求解}
\begin{lstlisting}[language=Python]
ampl = AMPL()
ampl.eval("... (先ほどのAMPLモデル文字列) ...")

# 1. 集合データの投入
ampl.set["CUSTOMERS"] = [c["id"] for c in consumers]
ampl.set["FACILITIES"] = [f["id"] for f in facilities]

# 2. パラメータデータの投入（辞書形式で直接注入可能）
ampl.param["w_setup"] = 2
ampl.param["w_dist"] = 6
ampl.param["demand"] = {c["id"]: c["demand"] for c in consumers}
ampl.param["capacity"] = {f["id"]: f["capacity"] for f in facilities}
ampl.param["setup_cost"] = {f["id"]: f["setupCost"] for f in facilities}
ampl.param["dist"] = dist_matrix  # タプルキー (i, j) の辞書

# 3. ソルバーに HiGHS を指定して求解実行
ampl.option["solver"] = "highs"
ampl.solve()
\end{lstlisting}
\end{frame}

\begin{frame}[fragile]{Python 実装 (3)：結果の抽出と集計}
\begin{lstlisting}[language=Python]
# 目的関数値（加重コスト）の取得
obj_val = ampl.get_objective("TotalCost").value()
ampl_soft_score = -int(round(obj_val))

# 開設された施設（y_j > 0.5 となる施設ID）を取得
y_sol = ampl.get_variable("y").get_values().to_dict()
ampl_open_facs = [j for j, val in y_sol.items() if val > 0.5]

# 顧客の割当先施設（x_ij > 0.5 となるペア）を取得
x_sol = ampl.get_variable("x").get_values().to_dict()
ampl_assignment = {i: j for (i, j), val in x_sol.items() if val > 0.5}

# 施設開設コスト合計と実移動距離合計を算出
total_setup = sum(setup_cost[j] for j in ampl_open_facs)
total_dist = sum(dist_matrix[(i, ampl_assignment[i])] for i in c_ids)
\end{lstlisting}
\end{frame}

\section{最適化結果と Timefold 解の比較}

\begin{frame}{最適化結果の比較サマリー}
  \begin{table}\footnotesize
    \centering
    \setlength{\tabcolsep}{6pt}
    \begin{tabular}{lrrr}
      \toprule
      \textbf{指標} & \textbf{JSON (Timefold)} & \textbf{AMPL (最適解)} & \textbf{改善効果} \\
      \midrule
      求解ステータス & メタヒューリスティクス & \alert{\textbf{大域的最適解 (Global)}} & 最適性保証 \\
      Soft Score & $-1{,}510{,}342$ & $\mathbf{-1{,}467{,}636}$ & \alert{\textbf{+42,706 改善}} \\
      総ペナルティ（目的関数） & $1{,}510{,}342$ & $\mathbf{1{,}467{,}636}$ & \alert{\textbf{2.83\% 削減}} \\
      \midrule
      開設施設数 & 7 施設 & 7 施設 & 同数 \\
      施設開設コスト合計 & 313,178 & \textbf{311,028} & \textbf{2,150 削減} \\
      顧客移動距離合計 (m) & 147,331 & \textbf{140,930} & \textbf{6,401 m 短縮} \\
      \midrule
      開設コスト寄与 ($w=2$) & 626,356 & \textbf{622,056} & $-4,300$ \\
      距離コスト寄与 ($w=6$) & 883,986 & \textbf{845,580} & $-38,406$ \\
      \bottomrule
    \end{tabular}
  \end{table}
  \vspace{-0.6em}
  \scriptsize
  ※ 開設施設数はどちらも 7 施設ですが、AMPL 最適解は開設費用と移動距離の\alert{\textbf{両方を同時に削減}}しています。
\end{frame}

\begin{frame}{開設施設と顧客割当の差異分析}
  \begin{columns}[T]
    \begin{column}{0.48\textwidth}
      \textbf{開設施設の見直し}
      \begin{itemize}\footnotesize
        \item \textbf{共通開設 (5施設)}: \\
          \{10, 11, 15, 22, 25\}
        \item \textbf{Timefold のみ開設 (閉鎖)}: \\
          \alert{\textbf{\{13, 26\}}} (開設費: 46,095 / 47,846)
        \item \textbf{AMPL 最適解で新設}: \\
          \alert{\textbf{\{19, 24\}}} (開設費: 48,631 / 43,160)
      \end{itemize}
      \vspace{0.5em}
      \textbf{顧客割当の変更}
      \begin{itemize}\footnotesize
        \item 全60顧客中 \alert{\textbf{18顧客 (30\%)}} の割当先施設が変更。
        \item より近接した施設への最適再配分が実現。
      \end{itemize}
    \end{column}
    \begin{column}{0.5\textwidth}
      \textbf{施設需要の充足量（上限容量 150）}
      \begin{table}\tiny
        \centering
        \begin{tabular}{crrr}
          \toprule
          施設ID & 開設費用 & JSON需要 & AMPL需要 \\
          \midrule
          10 & 42,682 & 90 & 105 \\
          11 & 46,001 & 135 & 135 \\
          13 & 46,095 & 150 & \alert{\textbf{0 (閉鎖)}} \\
          15 & 43,421 & 135 & \textbf{150 (満床)} \\
          19 & 48,631 & 0 & \alert{\textbf{120 (新設)}} \\
          22 & 33,084 & 135 & \textbf{150 (満床)} \\
          24 & 43,160 & 0 & \alert{\textbf{90 (新設)}} \\
          25 & 54,049 & 150 & \textbf{150 (満床)} \\
          26 & 47,846 & 105 & \alert{\textbf{0 (閉鎖)}} \\
          \bottomrule
        \end{tabular}
      \end{table}
    \end{column}
  \end{columns}
\end{frame}

\section{まとめ}

\begin{frame}{手法の比較とまとめ}
  \begin{itemize}\small
    \item \textbf{Timefold (局所探索・メタヒューリスティクス)}:
      \begin{itemize}\footnotesize
        \item 複雑なビジネスルールや動的変更に柔軟に対応できる。
        \item 本問題でも大域的最適解からわずか約 2.8\% 差の実用上良好な解を瞬時に探索。
        \item ただし、解の「真の最適性」は保証されない。
      \end{itemize}
    \item \textbf{AMPL + HiGHS (0-1整数線形計画法 / 厳密解法)}:
      \begin{itemize}\footnotesize
        \item 目的関数・制約条件を明示的に定式化することで、\alert{\textbf{厳密な大域的最適解}}を算出。
        \item 強定式化（Linking Constraint）により、HiGHS でわずか 1 分枝ノード（探索瞬時）で最適解を発見。
        \item 開設費用を抑えつつ距離を 6.4 km 短縮し、42,706 ポイントのスコア改善を達成。
      \end{itemize}
    \item \textbf{結論}:
      \begin{itemize}\footnotesize
        \item 実務では、厳密解法（AMPL）でベースラインや大域的最適解を算定し、ヒューリスティクスの評価・改善指標とするアプローチが極めて有効である。
      \end{itemize}
  \end{itemize}
\end{frame}

\end{document}
